% !TEX program = xelatex
% 编译方式：在本目录执行两遍 xelatex（详见 docs/_manual_common/README.md）。
% 章节源文件位于 chapters/，由本文件 \input 引入。
% 源码依据：include/hacdcpf/network_reconfiguration/、
%           src/network_reconfiguration/、tests/run_gui_server.cpp 与注册测试。
\documentclass[UTF8,fontset=fandol,a4paper,11pt]{ctexart}
\usepackage{../../_manual_common/hysim_manual}

\title{\textbf{交直流混合高弹性能源电力系统仿真平台}\\[0.5em]
       {\Large 网络重构技术手册}\\[0.3em]
       {\large 混合 AC/DC 拓扑 MILP、设备动作、兼容 ONR 与运行时验证}}
\author{HySim-XJTU-HRPES（西安交通大学 HRPES 团队）}
\date{\today}

\begin{document}
\maketitle
\begin{abstract}
\noindent 本手册以当前可达代码路径为边界，系统说明
\srcpath{run\_topology\_reconfiguration} 的混合 AC/DC 单时刻重构 MILP、
\srcpath{solve\_optimal\_reconfiguration} 的历史 ONR 兼容接口，以及
\srcpath{POST /api/session/run\_reconfig} 在核心候选解之上追加的富模型回投影、潮流与 OPF
验证。内容覆盖规范投影、AC/DC 域限定索引、源与负荷聚合、统一/分域辐射约束、可选
LinDistFlow、开关设备能力与保护联锁、启发式和 MILP 后端、求解证书、结果字段、HTTP 契约、
注册测试与审计发现。

\smallskip
\noindent\textbf{核心口径}：核心 C++ 重构函数返回的是图模型或 LinDistFlow MILP 的候选拓扑，
不会自行执行完整混合 PF/OPF；\fld{reconf\_loss\_mw} 是名义 1 pu 电流下的电阻拓扑代理，
不是求解潮流损耗；\fld{milp\_objective} 是多项加权目标且可含省略常数，不能换算成 MW。
完整 PF/OPF 和可执行性标志只在当前 HTTP 路由的后处理层赋值。

\smallskip
\noindent\textbf{文档层次}：第 2 章为工程参考级通用理论，与第 4 章的源码等价模型分立；
理论中的严格 DistFlow、二阶锥松弛、显式可靠性时序和多时段恢复均不自动构成当前实现承诺。
\end{abstract}

\tableofcontents
\newpage

% overview_contract.tex —— 网络重构模块总览与公共契约

\section{模块总览与公共契约}\label{sec:nr-overview}

\subsection{职责与非职责}

网络重构模块解决给定单时刻设备状态下的离散拓扑选择问题。核心调用链为
\begin{equation}
\mathcal S_{rich}
\xrightarrow{\Pi_{canon}}
\widehat{\mathcal S}
\xrightarrow{\text{topology/LinDistFlow MILP}}
\widehat z^star
\xrightarrow{\Phi^{-1}}
\text{rich-device actions}.
\end{equation}
其中 $\mathcal S_{rich}$ 是 \texttt{HybridPowerSystem}，$\widehat z^star$ 是规范边空间中的
开闭状态，$\Phi^{-1}$ 通过 \texttt{BranchExpandMap} 将变化边归因回开关、断路器或物理支路。

本模块的职责包括：
\begin{itemize}
  \item AC、DC 和 VSC 边的域限定故障/候选标识；
  \item 统一树或分域辐射森林约束、可选 DC 网状运行；
  \item 可选有功/无功平衡、电压降、箱式热限和切负荷；
  \item 富模型开断能力、闭锁和上游保护联锁过滤；
  \item 图启发式、外部 MILP 后端、原生 B\&C 显式路径及后验可行性检查；
  \item 规范边到设备动作的回投影和确定性动作排序。
\end{itemize}

核心 C++ 函数\textbf{不负责}完整非线性交直流潮流、完整混合 OPF、保护暂态、时序恢复、
储能 SOC 演化或可靠性指标积分。这些能力分别属于 power\_flow、optimal\_power\_flow、
reliability/resilience 等模块。生产 HTTP 路由在核心候选解之后额外调用 PF/OPF，但该后处理
不是 \srcpath{run\_topology\_reconfiguration} 本身的行为。

\subsection{代码所有权与可达路径}

\begin{center}\small
\begin{tabular}{@{}L{5.2cm}L{4.0cm}L{5.4cm}@{}}
\toprule
符号/入口 & 所有权 & 当前可达行为\\
\midrule
\srcpath{run\_topology\_reconfiguration} &
\srcpath{topology\_reconfiguration.cpp} &
混合 AC/DC 规范投影、MILP/启发式、结果提取、设备动作\\
\srcpath{solve\_optimal\_reconfiguration(ACSystem)} &
\srcpath{topology\_analysis.cpp} &
包装为 \texttt{HybridPowerSystem} 后调用上一入口，再运行验证 PF\\
\srcpath{solve\_optimal\_reconfiguration} 的混合重载 &
\srcpath{topology\_analysis.cpp} &
先规范投影，再仅把规范 AC 子系统传给 AC 兼容入口\\
\srcpath{is\_connected/is\_radial/count\_islands} &
\srcpath{topology\_analysis.cpp} &
只检查 \texttt{ACSystem} 中在运 \texttt{ACBranch}\\
\srcpath{POST /api/session/run\_reconfig} &
\srcpath{tests/run\_gui\_server.cpp} &
核心求解后回写富模型并追加拓扑、PF、AC OPF 与 JSON 响应\\
\bottomrule
\end{tabular}
\end{center}

\begin{quote}\small\textbf{重要：}
\srcpath{topology\_analysis.cpp} 在兼容包装返回之后仍保留一段旧 AC LinDistFlow B\&C 实现，
但当前控制流在该代码之前无条件 \texttt{return result}。因此旧段落不是运行时模型，头文件中
“模块自有 B\&C + MIR”描述也不能作为现行行为证据。
\end{quote}

\subsection{命名空间、头文件与门面边界}

公共类型均位于 \texttt{hacdcpf::analysis}：
\begin{itemize}
  \item \srcpath{topology\_reconfiguration.hpp}：
        \texttt{BranchRef}、\texttt{TopoReconfOptions}、\texttt{TopoReconfResult}；
  \item \srcpath{topology\_analysis.hpp}：
        三个 AC 拓扑查询、\texttt{ONROptions}、\texttt{ONRResult} 与兼容入口。
\end{itemize}
顶层 \texttt{hacdcpf} 门面当前不重导出这些头；库调用方应直接包含对应模块头文件。

\subsection{单位制、方向与索引空间}

\begin{center}\small
\begin{tabular}{@{}L{3.6cm}L{3.0cm}L{8.0cm}@{}}
\toprule
对象 & 单位/空间 & 契约\\
\midrule
组件 \fld{index} & 稳定 ID & 对外标识；不得当作 vector position\\
规范母线位置 & 0-based & AC 为 $[0,n_{ac})$，DC 为 $[n_{ac},n_b)$；只在函数内部使用\\
\texttt{BranchRef} & 类型化稳定 ID & \fld{category}+\fld{index}；混合系统首选\\
\fld{open\_branch\_ids} 等 & 裸整数 & 兼容字段；AC/DC/VSC 同号时有歧义\\
$P,Q,P_d,Q_d$ & p.u.（模型内） & 由 MW/Mvar 除以 \fld{base\_mva}；结果切负荷再乘回 MW/Mvar\\
$v_i$ & pu$^2$ & LinDistFlow 平方电压变量\\
$F_e,F_g$ & 无量纲 & 单商品虚拟流，仅证明拓扑连通性\\
\fld{max\_time\_s} & s & MILP 墙钟时限\\
\fld{reconf\_loss\_mw} & 名义 MW 代理 & $S_B\sum_{e\in E_{line}}r_e\beta_e$，不是 PF 损耗\\
\bottomrule
\end{tabular}
\end{center}

统一边位置固定为 AC 支路、DC 支路、VSC：
\begin{equation}
E=[E_{ac},E_{dc},E_{vsc}],\qquad
e_{dc}=n_{l,ac}+k,\quad e_{vsc}=n_l+k.
\end{equation}
该位置布局只用于矩阵装配；结果必须恢复为 \texttt{BranchRef}。

\subsection{选项完整契约}

\subsubsection{故障与候选集}
\begin{paramtable}
\fld{line\_failures} & $\mathcal F_{ac}$ & 稳定 ID & AC-only & 空 & 强制断开的 \texttt{ACBranch::index}\\
\fld{faulted\_branches} & $\mathcal F$ & \texttt{BranchRef} & AC/DC/VSC & 空 & 混合系统故障集\\
\fld{switchable\_branches} & $\mathcal C$ & \texttt{BranchRef} & AC/DC/VSC & 空 & 显式候选集\\
\fld{switchable\_branch\_ids} & $\mathcal C_{ac}$ & 稳定 ID & AC-only & 空 & 旧式 AC 候选集\\
\end{paramtable}
若两个候选字段均为空，物理 AC 支路和 DC 支路仅把初始断开边自动视作 tie；来自富模型
Switch/CircuitBreaker 的规范 AC 边则按设备能力自动判断。VSC 只有被结构化显式列出时才可切换。
一旦任一显式候选字段非空，未列出的分支保持初始状态。

\subsubsection{模型组与拓扑模式}
\begin{paramtable}
\fld{enable\_pf} & -- & 布尔 & -- & \texttt{true} & 创建 $P,Q,P_g,Q_g,v,s^P,s^Q$\\
\fld{enable\_voltage} & G4 & 布尔 & -- & \texttt{true} & PF 开启时施加线路电压降\\
\fld{enable\_thermal} & G5 & 布尔 & -- & \texttt{true} & PF 开启时施加 P/Q 箱式容量\\
\fld{split\_domain\_trees} & -- & 布尔 & -- & \texttt{false} & AC/DC 分域树，VSC 不计边数\\
\fld{allow\_dc\_mesh} & -- & 布尔 & -- & \texttt{false} & 启用分域模式且不施加 DC 树边数\\
\fld{loss\_aware} & -- & 布尔 & -- & \texttt{false} & PF 开启时使用 $t\ge |P|$ 代理\\
\end{paramtable}
\fld{allow\_dc\_mesh} 隐含分域行为。它允许 DC 环网，因此返回结果不能被解释为“整个混合系统
均执行了辐射约束”。

\subsubsection{权重、边界与求解控制}
\begin{paramtable}
\fld{lambda\_switch} & $\lambda_{sw}$ & -- & 建议 $\ge0$ & $1$ & 每次设备动作成本权重\\
\fld{lambda\_loss} & $\lambda_{loss}$ & -- & 建议 $\ge0$ & $10$ & 电阻/潮流绝对值代理权重\\
\fld{lambda\_shed} & $\lambda_{shed}$ & -- & 建议 $\ge0$ & $10^4$ & P/Q 切负荷权重\\
\fld{lambda\_island} & $\lambda_{isl}$ & -- & 建议 $\ge0$ & $10^5$ & 非主根启用权重\\
\fld{max\_switch\_ops} & $N_{sw}$ & 次数成本 & $0$=不限制 & $0$ & 隔离设备一次变化按 3 计\\
\fld{v\_min\_pu} & $\underline V$ & pu & $>0$ & $0.95$ & 母线未给有效下限时使用\\
\fld{v\_max\_pu} & $\overline V$ & pu & $>0$ & $1.05$ & 母线未给有效上限时使用\\
\fld{default\_rate\_mva} & $\overline S_0$ & MVA & $>0$ & $10$ & 支路/VSC 容量回退\\
\fld{big\_m\_v} & $M_{cap}$ & pu$^2$ & $>0$ & $2$ & 支路自适应 M 的上限\\
\fld{max\_time\_s} & -- & s & $>0$ & $300$ & 求解器时限\\
\fld{mip\_gap} & $\epsilon_{mip}$ & -- & $\ge0$ & $0.01$ & 相对间隙阈值\\
\fld{skip\_heuristic} & -- & 布尔 & -- & \texttt{false} & 强制进入 MILP 路径\\
\fld{solver} & -- & 字符串 & 见第~\ref{sec:nr-solver}~章 & \texttt{auto} & 后端选择\\
\fld{verbose} & -- & 布尔 & -- & \texttt{false} & spdlog 诊断\\
\end{paramtable}

公共入口没有集中校验上述数值域；负权重、反向电压区间、非正容量或未知求解器字符串的行为
不应视作稳定输入契约。生产调用方应先做输入校验。

\subsection{结果族速览}

\texttt{TopoReconfResult} 同时包含四类信息：
\begin{enumerate}
  \item 拓扑：\fld{open\_branches}/\fld{closed\_branches}、\fld{switched\_on}/
        \fld{switched\_off}；
  \item 富设备动作：\fld{switch\_operations}、\fld{ordered\_switch\_actions} 和序列状态；
  \item 优化：\fld{milp\_objective}、\fld{obj\_terms}、切负荷、求解器状态/间隙；
  \item 作用域：\fld{model\_scope} 与 \texttt{ValidityFlags}。
\end{enumerate}
\fld{feasible} 只表示返回向量通过当前线性模型的后验残差、整数性和边界检查；
\fld{optimal} 与 \fld{proven\_optimal} 只在后端状态被识别为最优且间隙满足阈值时为真。
图启发式成功时两者均为假。

\texttt{ONRResult} 是兼容形状，只保留 AC 整数支路 ID、代理目标、验证 PF 和
\texttt{BCStats}。它没有 \texttt{model\_scope}、限制列表、后端名称、切负荷和可执行性字段，
不能承载混合核心的全部证书。

% theory_reconfiguration.tex —— 网络重构通用理论
% 本文件为理论层章节，遵循 docs/modules/THEORY_WRITING_GUIDE.md。

\section{网络重构通用理论}\label{sec:nr-theory}

\begin{theorynote}
本章为工程参考级通用数学理论，覆盖图上的配电网重构、辐射森林、单商品流连通性、
DistFlow/LinDistFlow、开关动作、混合 AC/DC 分域拓扑、目标标度和验证层次。公式不要求
逐式对应代码；与平台实现的对应关系见章末。严格支路流、二阶锥松弛、多时段恢复和
N--1 安全约束若超出当前实现，均以“未实现扩展说明”标记。
\end{theorynote}

\subsection{记号与问题定义}

令电气网络为图 $G=(V,E)$，候选边集 $E=E_f\cup E_s$，其中 $E_f$ 是状态固定边，
$E_s$ 是可操作开关边。$z_e\in\{0,1\}$ 表示边 $e$ 最终闭合，$z_e^0$ 是初始状态，
$a_e\in\{0,1\}$ 表示故障可用性。源集合为 $R\subseteq V$，负荷为 $(p_i^d,q_i^d)$。
单时刻网络重构的一般形式为
\begin{equation}
\min_{z,x,s}\quad C_{sw}(z,z^0)+C_{loss}(x,z)+C_{shed}(s)+C_{isl}(z)
\end{equation}
\begin{equation}
\text{s.t.}\quad z_e\le a_e,\quad
z_e=z_e^0\ (e\in E_f),\quad
(z,x,s)\in\mathcal T\cap\mathcal P,
\end{equation}
其中 $\mathcal T$ 是拓扑可行集，$\mathcal P$ 是选定电气近似的可行集。重构问题的困难来自
$z$ 的组合性，以及开边时电压降/潮流方程的析取结构。

\subsection{辐射树、森林与根节点}

\begin{definition}[辐射供电拓扑]
对连通分量 $G_k=(V_k,E_k)$，若 $|E_k|=|V_k|-1$ 且连通，则 $G_k$ 为树。允许多个有源或
失电分量时，最终拓扑为森林；若森林含 $r$ 个分量，则
\begin{equation}
\sum_{e\in E}z_e=|V|-r.
\end{equation}
\end{definition}

仅有边数等式不能排除“一个分量含环、另一个分量断开”的抵消，必须与连通性或分量约束联用。
当每个供电分量选择一个根 $y_r\in\{0,1\}$ 时，可写为
\begin{equation}
\sum_{e\in E}z_e+\sum_{r\in R}y_r=|V|,
\end{equation}
但该式仍需虚拟流保证每个非根节点可由某个已选根到达。

\begin{proposition}[单商品流的连通性证书]
给每条候选边有符号虚拟流 $f_e$，给根候选虚拟注入 $g_r\ge0$。取一致的关联矩阵 $A$，若
\begin{align}
A f-Bg&=-\mathbf 1,\\
|f_e|&\le |V|z_e,\\
0\le g_r&\le |V|y_r,
\end{align}
并满足树/森林基数关系，则每个节点位于某个已选根的树中。这里 $B$ 把根注入映射到节点。
\end{proposition}

直观上每个节点消费一单位虚拟商品，只有根能够注入；未闭合边不能传商品。因此任何无根分量
都会违反平衡。流方向与电功率方向无关，只要关联矩阵符号全局一致即可。

\subsection{故障、固定状态与开关预算}

故障与设备能力共同决定状态界：
\begin{equation}
\begin{array}{ll}
a_e=0:&z_e=0,\\
e\notin E_s:&z_e=z_e^0,\\
e\in E_s,\ a_e=1:&0\le z_e\le1,\ z_e\in\mathbb Z.
\end{array}
\end{equation}
若动作成本为 $c_e^{op}>0$，总预算为
\begin{equation}
\sum_{e\in E_s}c_e^{op}|z_e-z_e^0|\le N_{op}.
\end{equation}
对二元 $z_e$，绝对值可以按初始状态展开为 $z_e$ 或 $1-z_e$，无需新增变量。

物理开关动作不等于单个二元状态变化。隔离开关可能要求“上游保护断开--隔离设备操作--上游
重合”三步序列，因此操作成本应来自设备能力，而非只数拓扑边。

\subsection{DistFlow 与 LinDistFlow}

对从 $i$ 指向 $j$ 的径向 AC 支路，经典 branch-flow/DistFlow 方程可写为
\begin{align}
P_{ij}&=p_j^d-p_j^g+\sum_{k:j\to k}P_{jk}+r_{ij}\ell_{ij},\\
Q_{ij}&=q_j^d-q_j^g+\sum_{k:j\to k}Q_{jk}+x_{ij}\ell_{ij},\\
v_j&=v_i-2(r_{ij}P_{ij}+x_{ij}Q_{ij})+(r_{ij}^2+x_{ij}^2)\ell_{ij},\\
P_{ij}^2+Q_{ij}^2&=v_i\ell_{ij},
\end{align}
其中 $v_i=|V_i|^2$，$\ell_{ij}=|I_{ij}|^2$。忽略二阶损耗项并固定方向后得到
LinDistFlow：
\begin{align}
P_{ij}&\approx p_j^d-p_j^g+\sum_{k:j\to k}P_{jk},\\
Q_{ij}&\approx q_j^d-q_j^g+\sum_{k:j\to k}Q_{jk},\\
v_j-v_i+2(r_{ij}P_{ij}+x_{ij}Q_{ij})&\approx0.
\end{align}

当支路可开断时，电压式用 Big-M 析取：
\begin{equation}
-M_e(1-z_e)\le v_j-v_i+2(r_eP_e+x_eQ_e)\le M_e(1-z_e).
\end{equation}
容量可以采用箱式近似 $|P_e|,|Q_e|\le\bar S_ez_e$，但它既允许角点
$(\bar S_e,\bar S_e)$，也比圆形 $P_e^2+Q_e^2\le\bar S_e^2z_e$ 更宽松。

\begin{gapnote}
当前模块不创建 $\ell_e$，不保留 $r_e\ell_e/x_e\ell_e$ 损耗项，也不施加
$P_e^2+Q_e^2\le v_i\ell_e$ 的二阶锥约束。因此结果不是严格 DistFlow、SOCP 下界或 AC 可行性
证书；完整电气可行性必须由后续非线性 PF/OPF 检查。
\end{gapnote}

\subsection{切负荷与孤岛可行性}

用 $s_i^P,s_i^Q\ge0$ 表示削减后，节点平衡可写为
\begin{align}
\sum_{e\in\delta^+(i)}P_e-\sum_{e\in\delta^-(i)}P_e-P_i^g-s_i^P&=-P_i^d,\\
\sum_{e\in\delta^+(i)}Q_e-\sum_{e\in\delta^-(i)}Q_e-Q_i^g-s_i^Q&=-Q_i^d.
\end{align}
零容量伪根可以证明一个失电分量的图连通性，但不能注入真实功率；若该分量保持孤立，平衡式迫使
负荷被削减。工程上通常还应施加 $0\le s_i^P\le P_i^d$ 及适当的无功削减边界。

\begin{gapnote}
当前核心实现只给 $s_i^P,s_i^Q$ 非负下界，没有按节点需求设置上界，也没有显式耦合恒功率因数
削减。对常规正负荷和高切负荷惩罚，该松弛通常取不超过需求的解；但这不是任意输入下的数学保证。
\end{gapnote}

\subsection{损耗目标与多目标标度}

严格有功损耗为 $\sum_e r_e\ell_e$。常见近似包括：
\begin{align}
L_{status}(z)&=\sum_e r_ez_e,\\
L_{abs}(P)&=\sum_e r_et_e,\qquad t_e\ge|P_e|,\\
L_{quad}(P,Q,v)&\approx\sum_e r_e\frac{P_e^2+Q_e^2}{v_i}.
\end{align}
$L_{status}$ 只偏好低电阻闭边组合；$L_{abs}$ 随流量变化但不是 $I^2R$；$L_{quad}$ 更接近物理，
却需要二次/锥建模。多目标加权
\begin{equation}
J=\lambda_{sw}C_{sw}+\lambda_{loss}L+\lambda_{shed}S+\lambda_{isl}I
\end{equation}
只有在各项尺度和优先级经过校准时才有稳定语义。若业务要求“先最小切负荷，再最小动作，再最小
损耗”，字典序或分层求解比仅凭大权重更可靠。

\begin{gapnote}
当前模块采用单次加权和，不执行字典序优化，也不报告目标各项的归一化尺度。默认大惩罚只是工程
偏好，不构成严格优先级证明。
\end{gapnote}

\subsection{混合 AC/DC 的分域拓扑}

混合图含 $G_{ac}=(V_{ac},E_{ac})$、$G_{dc}=(V_{dc},E_{dc})$ 和换流边 $E_c$。配电辐射性通常
针对每个电气域定义：
\begin{equation}
\sum_{e\in E_{ac,k}}z_e=|V_{ac,k}|-1,
\qquad
\sum_{e\in E_{dc,k}}z_e=|V_{dc,k}|-1,
\end{equation}
换流器作为跨域功率端口，不计入域内树边数。若允许 MTDC 网状运行，可去掉 DC 边数等式但仍保留
电气平衡、容量和参考条件。

统一混合树是另一种合法建模选择：把换流器当作普通图边，对 $V_{ac}\cup V_{dc}$ 施加一棵树。
它比“AC 径向 + DC 可网状 + 换流器自由”的业务语义更严格，不能把两者混称为同一径向约束。

\subsection{候选解、最优证书与可执行方案}

三个层次必须分开：
\begin{enumerate}
  \item \textbf{线性模型可行}：变量满足当前 MILP 的等式、不等式、边界和整数性；
  \item \textbf{已证明最优}：求解器状态和 MIP gap 对该模型提供最优性证书；
  \item \textbf{工程可执行}：动作能映射回设备，序列满足能力/保护约束，回投影后 PF/OPF 通过。
\end{enumerate}
第 1 层不蕴含第 2 或第 3 层。尤其是图启发式可提供第 1 层 incumbent，但不能提供最优性；
LinDistFlow 的最优解也可能在 AC PF 中不可行。

\begin{gapnote}
当前 C++ 核心只完整承担第 1 层，并在部分后端状态下标注第 2 层；第 3 层的完整组合只在生产 HTTP
路由后处理中计算。库调用方必须自行复现该验证链，不能仅检查 \fld{feasible}。
\end{gapnote}

\subsection{未实现扩展边界}

下列常见 ONR 扩展不属于当前模块：
\begin{itemize}
  \item 多时段开关计划、最小停留时间、储能 SOC、移动电源路由；
  \item N--1 安全重构、随机/鲁棒负荷与可再生不确定性；
  \item 三相不平衡、相别开关、零序/中性线与相间耦合；
  \item 保护动作时间、电弧/涌流、同步检查和故障电流约束；
  \item 精确 $I^2R$、SOCP/SDP 松弛最优性界和完整 AC/DC 非线性联合重构；
  \item 拓扑可观测性、通信可用性和遥控失败概率。
\end{itemize}
这些能力若由 reliability、resilience 或其他模块提供，也不能倒推为本入口的隐式功能。

\subsection{与实现的对应}

\begin{longtable}{@{}L{7.4cm}L{7.0cm}@{}}
\toprule
理论条目 & 实现状态\\
\midrule
\endfirsthead
\toprule
理论条目 & 实现状态\\
\midrule
\endhead
故障/固定/候选状态界 & \implfull{topology\_reconfiguration.cpp:run\_topology\_reconfiguration}\\
统一树的单商品流与根变量 & \implfull{topology\_reconfiguration.cpp:run\_topology\_reconfiguration}\\
AC/DC 分域树与 DC mesh 选项 & \implfull{topology\_reconfiguration.cpp:run\_topology\_reconfiguration}\\
LinDistFlow 节点平衡与电压降 & \implpart{无支路损耗项、无电流平方变量、VSC 无功为端口近似}\\
圆形/锥形视在功率容量 & \implpart{仅实现 P/Q 独立箱约束}\\
动作预算与设备动作成本 & \implfull{topology\_reconfiguration.cpp:run\_topology\_reconfiguration}\\
严格 $I^2R$ 或 SOCP 损耗 & \implnone\\
多目标字典序求解 & \implnone\\
后验线性残差/整数性/边界证书 & \implfull{topology\_reconfiguration.cpp:run\_topology\_reconfiguration}\\
完整 PF/OPF 工程可执行性 & \implpart{tests/run\_gui\_server.cpp:POST /api/session/run\_reconfig；核心函数不执行}\\
多时段、不确定性、N--1 与三相 ONR & \implnone\\
\bottomrule
\end{longtable}

\subsection{参考文献}

\begin{thebibliography}{9}\small
\bibitem{nr:baranwu}
M. E. Baran and F. F. Wu,
``Network reconfiguration in distribution systems for loss reduction and load balancing,''
\emph{IEEE Transactions on Power Delivery}, vol. 4, no. 2, pp. 1401--1407, 1989.

\bibitem{nr:civanlar}
S. Civanlar, J. J. Grainger, H. Yin, and S. S. H. Lee,
``Distribution feeder reconfiguration for loss reduction,''
\emph{IEEE Transactions on Power Delivery}, vol. 3, no. 3, pp. 1217--1223, 1988.

\bibitem{nr:lavorato}
M. Lavorato, J. F. Franco, M. J. Rider, and R. Romero,
``Imposing radiality constraints in distribution system optimization problems,''
\emph{IEEE Transactions on Power Systems}, vol. 27, no. 1, pp. 172--180, 2012.

\bibitem{nr:jabr}
R. A. Jabr, R. Singh, and B. C. Pal,
``Minimum loss network reconfiguration using mixed-integer convex programming,''
\emph{IEEE Transactions on Power Systems}, vol. 27, no. 2, pp. 1106--1115, 2012.
\end{thebibliography}

% canonical_device_semantics.tex —— 规范投影、边空间与设备动作

\section{规范投影与设备语义}\label{sec:nr-projection}

\subsection{投影配置与所有权}

核心入口首先调用
\srcpath{projection::RichToCanonicalOperator::apply}，固定配置为
\begin{center}
\texttt{strip\_dead\_islands=false},\qquad
\texttt{preserve\_switch\_branches=true}.
\end{center}
前者保留当前失电但可能经 tie 恢复的区段；后者确保富模型 Switch/CircuitBreaker 在规范 AC
支路表中保留可归因边。求解模型只读取投影后的 \texttt{ac}、\texttt{dc} 和
\texttt{vsc\_converters}，但设备能力判断仍回看原始 \texttt{sys.ac.switches} 与
\texttt{sys.ac.circuit\_breakers}。

\begin{quote}\small\textbf{重要：}
投影是重构入口的一部分。调用方不应先用默认投影删除 dead island，再把结果传入重构函数；
否则已删除的失电区段和联络路径无法由内部投影恢复。
\end{quote}

\subsection{AC/DC 域限定母线映射}

规范位置采用两个独立映射：
\begin{equation}
\mu_{ac}:\texttt{ACBus::index}\mapsto[0,n_{ac}),\qquad
\mu_{dc}:\texttt{DCBus::index}\mapsto[n_{ac},n_b).
\end{equation}
VSC 端点解析为 $(\mu_{ac}(bus_{ac}),\mu_{dc}(bus_{dc}))$。这允许 AC 母线 1 和 DC 母线 1
同时存在，不发生哈希覆盖。任何跨域调用若改用裸 \texttt{unordered\_map<int,int>} 都会破坏
该不变量。

端点 ID 无法解析的边保留默认位置 $-1$，随后不参与矩阵装配。入口没有单独返回“跳过了无效端点边”
的诊断；严格调用方应先运行系统 validation。

\subsection{类型化边引用}

\texttt{BranchRef} 定义为
\begin{equation}
\operatorname{ref}(e)=(\operatorname{category}(e),\operatorname{component\ index}(e)),
\end{equation}
类别只承诺 \texttt{AC\_Line}、\texttt{DC\_Line}、\texttt{VSC\_Coupling}。以下字段必须用该
类型解释：\fld{faulted\_branches}、\fld{switchable\_branches}、\fld{open\_branches}、
\fld{closed\_branches}、\fld{switched\_on} 和 \fld{switched\_off}。

兼容整数向量把三类稳定 ID 混放，注释中所谓区间并不是可靠的数值分区，因为每个组件表的
\fld{index} 都可独立取值。混合系统中必须忽略裸整数结果并使用结构化字段。

\subsection{候选设备判定}

\subsubsection{物理 AC/DC/VSC 边}

没有 \texttt{BranchExpandMap} 来源的 AC 物理支路：显式候选模式下只有被列出的 ID 可变；自动模式
下只有初始断开支路可闭合。DC 支路采用相同规则，但只读取结构化 DC 候选。VSC 在自动模式下固定，
只有结构化显式列出才可切换。

故障优先级高于候选：任何匹配 \fld{line\_failures}/\fld{faulted\_branches} 的边都强制
$\beta_e=F_e=0$。未匹配的故障 ID 被静默忽略；当前没有“请求故障数/实际匹配数”字段。

\subsubsection{富模型 Switch}

投影来源为 Switch 时，\srcpath{effective\_switch\_capabilities} 与设备字段共同决定可变性：
\begin{itemize}
  \item 不在运或 \texttt{Fuse}：不可作为即时恢复动作；
  \item 初始闭合且 \fld{locked\_closed}：不可断开；
  \item 初始断开且 \fld{locked\_open}：不可闭合；
  \item 要求失电操作：必须绑定在运、非熔断器、具备开断能力的上游保护开关；
  \item 闭合恢复只允许 Tie、CircuitBreaker、Recloser 或 LoadBreakSwitch，且能力声明允许；
  \item 直接可带负荷/故障电流操作的设备，状态变化成本为 1；要求失电操作者保守计 3。
\end{itemize}
显式列出一个设备规范边不能越过上述硬过滤。

\subsubsection{CircuitBreaker 与其他来源}

投影来源为 CircuitBreaker 时，当前硬条件仅检查富断路器记录存在且 \fld{in\_service}；
Switch 的锁定/能力细分规则没有完整复用于 CircuitBreaker POD。Transformer 来源永不成为重构决策。

\subsection{重叠边去重}

富开关设备可能与其控制的物理 AC 支路共享端点。规范投影为了归因同时保留两条边；若两者都进入
拓扑，会产生虚假并联并破坏径向边数。核心实现扫描所有 AC 边：若 Switch/CircuitBreaker 来源边与
非设备来源边端点相同，设备等效边设置 \texttt{edge\_participates=false}；自环设备边同样排除。

被排除边仍按初始状态进入结果的 open/closed 列表，但虚拟流和 LinDistFlow 变量固定为零，不能
改变拓扑。这是一种端点启发式去重，不是显式“设备串接物理支路”的联合状态方程。

\subsection{源模型与根资格}

核心把同一规范母线上的资源聚合为一组 $P_g,Q_g$ 上下界。当前来源如下：
\begin{center}\small
\begin{tabular}{@{}L{4.4cm}L{4.0cm}L{6.0cm}@{}}
\toprule
资源 & 注入能力 & 可形成拓扑根\\
\midrule
AC Generator & $P_{min}/P_{max},Q_{min}/Q_{max}$ & 仅 \fld{is\_slack} 或 slack 母线\\
AC StaticGenerator/Renewable/PV/Storage & 有效容量与 Q 范围 & 否\\
AC ExternalGrid & 短路容量或默认容量 & 是\\
DC StaticGenerator/Storage/PVArray & 有效有功容量 & 是\\
无 VSC 连接的在运 DC\_V 母线 & $\max(\overline S_0,S_B)$ & 是\\
普通 slack 母线回退 & $\max(\overline S_0,S_B)$ & 是\\
失电分量伪根 & 零 P/Q 容量 & 是，仅虚拟连通\\
\bottomrule
\end{tabular}
\end{center}

VSC 相连的 DC\_V 母线只被视为电压参考，不被当作独立能源；其有功必须经 VSC 到达。若系统没有
任何真实可成根的 AC/DC 源，入口立即返回 \fld{feasible=false}，状态字符串含
\texttt{no AC/DC source bus available}。

当前闭合边图中每个无真实根分量会加入一个零容量伪根。它只满足拓扑商品流；PF 开启时，孤立分量
仍须以 $s^P/s^Q$ 平衡负荷。\fld{lambda\_island} 对除主根外的 $\gamma_g$ 加惩罚。

\subsection{负荷聚合}

负荷按规范空间\textbf{相加}：
\begin{align}
P_{d,i}^{ac}&=ACBus.pd+\sum_{Load@i}Load.p,\\
Q_{d,i}^{ac}&=ACBus.qd+\sum_{Load@i}Load.q,\\
P_{d,i}^{dc}&=DCBus.pd+\sum_{DCLoad@i}P_{effective}.
\end{align}
这与部分模块“存在 Load 表时替代母线聚合字段”的策略不同。注册回归专门固定了 AC 母线需求与
Load 记录相加、DCBus 与 DCLoad 相加的语义。

投影前的 FlexibleLoad、AsymmetricLoad、Charger 等富组件如何进入规范负荷，由 canonical
projection 合同决定；本模块不重复直接遍历这些富表。

\subsection{设备动作回投影}

求解后每条规范边比较初始状态与 $\beta_e^\star>0.5$。AC 边通过 \texttt{BranchExpandMap} 恢复为
Switch 或 CircuitBreaker；无来源映射、DC 和 VSC 变化以 \texttt{DeviceKind::Branch} 返回。
\texttt{SwitchOperation.index} 的索引空间取决于 \fld{kind}：设备种类时为富设备稳定 ID，Branch
时为相应类别组件 ID。

动作顺序按 \fld{switch\_operations} 的确定性遍历生成：
\begin{itemize}
  \item 可直接开断设备：单步 \texttt{open}/\texttt{close}；
  \item 要求失电操作的 Switch：上游 \texttt{open}，本设备 \texttt{open/close}，若最终仍需上游
        闭合则追加 \texttt{reclose}；
  \item 映射缺失、Fuse 或缺少可操作上游保护：序列无效并停止追加。
\end{itemize}
该序列只验证静态能力和绑定，不模拟保护动作时间、同步条件、故障电流、触头状态或操作失败概率。

\subsection{核心与 HTTP 回投影的差异}

核心函数只生成动作及 \fld{switch\_sequence\_validated}。生产 HTTP 路由另外：
\begin{enumerate}
  \item 把故障和动作写入富模型副本；
  \item 从修改后的富模型重新执行 canonical projection；
  \item 计算 AC/DC/混合拓扑；
  \item 运行完整 PF 和 AC OPF；
  \item 联合设置 \fld{post\_action\_reprojection\_validated}、
        \fld{post\_power\_flow\_validated}、\fld{full\_hybrid\_opf\_validated} 和
        \fld{executable}。
\end{enumerate}
因此直接 C++ 调用看到这些四个字段为默认 false 是当前实现行为，不表示核心已经尝试并失败。

\section{当前可达重构模型}

本章只描述 \srcpath{src/network\_reconfiguration/topology\_reconfiguration.cpp:run\_topology\_reconfiguration} 的
\srcpath{run\_topology\_reconfiguration}。历史入口
\srcpath{solve\_optimal\_reconfiguration} 在
\srcpath{src/network\_reconfiguration/topology\_analysis.cpp:solve\_optimal\_reconfiguration} 已提前包装并返回该混合求解器结果；其后保留的旧 AC MILP
在当前控制流不可达，故不作为现行模型。

\subsection{变量布局和边界传播}

统一边序为 AC、DC、VSC。拓扑变量包括有符号虚拟流 $F_e\in[-n_b,n_b]$、
源虚拟注入 $F_g\in[0,n_b]$、支路状态 $\beta_e$ 和根指示 $\gamma_g$。
启用 PF 时再创建线路/VSC 的 $P,Q$、源 $P_g,Q_g$、平方电压 $v_i$、有功切负荷
$s_i^P\ge0$、AC 无功切负荷 $s_i^Q\ge0$；仅在同时启用 loss aware 时创建
$t_e\ge0$。准确偏移见 \srcpath{src/network\_reconfiguration/topology\_reconfiguration.cpp:VarLayout}。

状态边界按以下分支执行：不参与拓扑的 overlay 边固定为原状态且虚拟流固定零；故障边
$\zeta_e<0.5$ 固定 $\beta_e=F_e=0$；合法可切换边保留 $0\le\beta_e\le1$ 并成为二元变量；
其余边固定为原状态，固定断开时同时固定 $F_e=0$。
同一当前连通分量内部的断开 tie 即便此前自由，也会固定断开。
依据为 \srcpath{src/network\_reconfiguration/topology\_reconfiguration.cpp:add\_dc\_source（lambda）}。

\subsection{实际目标函数}

对初始断开边，$\beta_e$ 的动作系数为 $+\lambda_{sw}c_e$；对初始闭合边为
$-\lambda_{sw}c_e$，被省略的常数是 $\lambda_{sw}c_e$。因此目标中的动作部分等价于
对状态变化计费，但存入 LP 的不是 $\lambda_{sw}\sum\beta_e$。
损耗分支为
\begin{equation}
 L_e=\begin{cases}
 \lambda_{loss}|r_e|t_e,&\fld{loss\_aware}\land\fld{enable\_pf},\\
 \lambda_{loss}|r_e|\beta_e,&\text{否则}.
 \end{cases}
\end{equation}
启用 PF 时再加 $\lambda_{shed}(\sum_i s_i^P+\sum_{i\in AC}s_i^Q)$；
除 primary root 外的每个 $\gamma_g$ 加 $\lambda_{island}\gamma_g$。
完整赋值位于 \srcpath{src/network\_reconfiguration/topology\_reconfiguration.cpp:log\_domain\_cardinality（lambda）}。

\subsection{拓扑等式与不等式}

统一树模式下，虚拟流平衡的矩阵符号为：from 端 $+F_e$、to 端 $-F_e$、源端
$-F_g$，每个母线右端为 $-1$；即
\begin{equation}
 \sum_{e:f(e)=i}F_e-\sum_{e:t(e)=i}F_e-\sum_{g:i(g)=i}F_g=-1.
\end{equation}
分域树模式把所有 VSC 虚拟流固定为零，并分别在每个 AC/DC 连通分量代表节点把右端
从 $-1$ 增加该分量母线数；见 \srcpath{src/network\_reconfiguration/topology\_reconfiguration.cpp:log\_domain\_cardinality（lambda）}。

统一模式的森林基数等式为
\begin{equation}
 \sum_{e\in\mathcal E_{participating}}\beta_e+\sum_g\gamma_g=n_b.
\end{equation}
分域模式不使用该等式，而对每个预计算 AC 分量固定
$\sum_{e\in component}\beta_e=|V_{component}|-1$；仅当不允许 DC mesh 时，对每个 DC 分量
也创建同式。VSC 不计入分域树边数。依据为
\srcpath{src/network\_reconfiguration/topology\_reconfiguration.cpp:log\_domain\_cardinality（lambda）}。

仅对自由 $\beta_e$ 创建
\begin{equation}
 -n_b\beta_e\le F_e\le n_b\beta_e.
\end{equation}
每个源创建 $F_g\le n_b\gamma_g$，并创建 $\sum_g\gamma_g\ge1$。
统一模式存在真实 DC 源时，另创建 DC 源集合上的 $\sum\gamma_g\ge1$。
依据为 \srcpath{src/network\_reconfiguration/topology\_reconfiguration.cpp:log\_domain\_cardinality（lambda）} 和
\srcpath{src/network\_reconfiguration/topology\_reconfiguration.cpp:bigm\_branch（lambda）}。

\subsection{可选 LinDistFlow 约束}

启用 PF 时，有功矩阵行采用 from $+P_e$、to $-P_e$、源 $-P_g$、切负荷 $-s_i^P$，
右端为 $-P_{d,i}$；AC 无功采用相同符号，右端 $-Q_{d,i}$。
这对应
\begin{align}
 \sum_{out}P_e-\sum_{in}P_e-\sum_gP_g-s_i^P&=-P_{d,i},\\
 \sum_{out}Q_e-\sum_{in}Q_e-\sum_gQ_g-s_i^Q&=-Q_{d,i}\quad(i\in AC).
\end{align}
VSC $P$ 同时进入 AC/DC 两端；VSC $Q$ 只进入属于 AC 的端点。
依据为 \srcpath{src/network\_reconfiguration/topology\_reconfiguration.cpp:log\_domain\_cardinality（lambda）}。

启用电压约束时，对 AC/DC 线路而非 VSC 创建
\begin{equation}
 |v_j-v_i+2r_eP_e+2x_eQ_e|\le M_e(1-\beta_e),
\end{equation}
DC 边没有 $Q$ 项。源码的支路大 M 为
\begin{equation}
 M_e=\min\!\left(\fld{big\_m\_v},
 \max\left(0.1,|\Delta v_{span}|+2(|r_e|+|x_e|)\bar S_e\right)\right),
\end{equation}
其中 $\Delta v_{span}$ 取两种端点上下界差的较大者；见
\srcpath{src/network\_reconfiguration/topology\_reconfiguration.cpp:log\_domain\_cardinality（lambda）}。

启用热限时，线路/VSC 有功均使用独立箱约束
$|P_e|\le\bar S_e\beta_e$；AC 线路与 VSC 无功使用
$|Q_e|\le\bar S_e\beta_e$。这不是圆形视在功率约束；DC 线路无 $Q$ 变量。
loss aware 分支创建 $P_e-t_e\le0$ 和 $-P_e-t_e\le0$，故 $t_e\ge|P_e|$；
注释称其为 $I^2R$ 代理，但实际目标是 $|r_e||P_e|$ 的线性代理。
依据为 \srcpath{src/network\_reconfiguration/topology\_reconfiguration.cpp:bigm\_branch（lambda）} 和
\srcpath{src/network\_reconfiguration/topology\_reconfiguration.cpp:bigm\_branch（lambda）}。

\subsection{开关预算与结果语义}

预算只计自由边，并乘设备动作成本 $c_e$：初始断开边贡献 $c_e\beta_e$，初始闭合边贡献
$c_e(1-\beta_e)$。矩阵形式把闭合边常数移到右端：
\begin{equation}
 \sum_{e:\alpha_e=0}c_e\beta_e-
 \sum_{e:\alpha_e=1}c_e\beta_e
 \le N_{max}-\sum_{e:\alpha_e=1}c_e.
\end{equation}
仅当 \fld{max\_switch\_ops}>0 时创建，见
\srcpath{topology\_reconfiguration.cpp:run\_topology\_reconfiguration}。

旧手册给出的 $\sum r_eS_B$ 不是当前 loss aware 目标。结果的
\fld{feasible}/\fld{optimal}/\fld{proven\_optimal} 还受后端状态字符串规范化和 primal 可用性影响；
详细提取位于 \srcpath{topology\_reconfiguration.cpp:run\_topology\_reconfiguration}。

\subsection{覆盖边界}

本章覆盖当前可达 MILP 的变量、目标和带源码标签的 T/P/G 约束。源容量聚合、设备开断能力、
overlay 去重、伪根生成、BFS 初解、后处理开关序列均是离散工程分支，源码区间分别为
\srcpath{src/network\_reconfiguration/topology\_reconfiguration.cpp:run\_topology\_reconfiguration}、
\srcpath{src/network\_reconfiguration/topology\_reconfiguration.cpp:bigm\_branch（lambda）}；本章不以通用图论公式替代这些实现。

% solver_result_contract.tex —— 求解路径、后验证与结果字段

\section{求解路径与结果证书}\label{sec:nr-solver}

\subsection{模型装配与二元变量}

矩阵以 Eigen 稀疏 triplet 装配为 \texttt{MIPModel}。只有自由 $\beta_e$ 和全部 $\gamma_g$ 被声明为
二元变量；固定 $\beta_e$ 保持连续类型但满足 $lb=ub\in\{0,1\}$，以减少 B\&C 分支规模。
故障边、不可操作边、排除 overlay 和同分量无效 tie 均通过变量边界固定。

\fld{enable\_pf=false} 时只创建 $F,F_g,\beta,\gamma$；此时
\fld{enable\_voltage}、\fld{enable\_thermal} 和 \fld{loss\_aware} 不创建相应变量或约束。
\fld{model\_scope} 分别固定为：
\begin{equation}
\begin{cases}
\texttt{hybrid-acdc-topology-lindistflow},&\fld{enable\_pf}=true,\\
\texttt{hybrid-acdc-topology-connectivity},&\fld{enable\_pf}=false.
\end{cases}
\end{equation}

\subsection{图启发式}

当 \fld{skip\_heuristic=false} 且不是分域模式时，核心先从主根沿当前健康闭合边做 BFS。若所有母线
已可达，则取 BFS 树；若有不可达节点，则从一条连接“已达/未达”两侧的初始断开自由边中选择当前
线性系数最小者，再扩展一次 BFS。启发式向量设置：
\begin{itemize}
  \item participating 边只闭合 BFS parent edge，overlay 保留原状态；
  \item 主根 $\gamma=1$，虚拟流由子树大小生成；
  \item 其余变量初始为零。
\end{itemize}
只有该向量同时满足全部等式、不等式和变量界（最大违反 $<10^{-6}$）时才被接受。PF 开启且存在
非零负荷时，零 $P/Q/P_g/s$ 往往使启发式失败，从而进入 MILP。

成功的启发式返回
\begin{center}
\texttt{solver\_backend=graph-heuristic},\quad
\texttt{solver\_status=feasible incumbent},\quad
\texttt{solver\_mip\_gap=inf},\quad
\texttt{proven\_optimal=false}.
\end{center}
它没有与 MILP 下界比较，不能称为最优 ONR。

\subsection{实际后端分派}

\begin{center}\small
\begin{tabular}{@{}L{2.5cm}L{6.0cm}L{6.0cm}@{}}
\toprule
\fld{solver} & 首选路径 & 失败行为\\
\midrule
\texttt{native} & \texttt{solve\_milp\_bc}，MIR、5 轮根割、伪成本分支 & 直接失败返回\\
\texttt{scip} & \texttt{ScipAdapter} & 不可用或失败则返回\\
\texttt{highs} & \texttt{StrictHighsBranchAndCutAdapter} & 再尝试 SCIP；不尝试 Native\\
\texttt{auto} & HiGHS & 再尝试 SCIP；不尝试 Native\\
其他字符串 & 与 \texttt{auto} 相同 & HiGHS 后 SCIP\\
\bottomrule
\end{tabular}
\end{center}

这一分派以 \srcpath{topology\_reconfiguration.cpp:run\_topology\_reconfiguration} 为准。
公共头注释中 \texttt{auto (HiGHS->native)} 与实现不符；当前实际没有 auto-to-native fallback。
外部后端成功时 \fld{bc\_stats} 保持默认值，诊断应读取 \fld{solver\_backend}、
\fld{solver\_status} 和 \fld{solver\_mip\_gap}。

\subsection{成功、可行与最优判定}

外部或 Native 返回 \texttt{stats.success=true} 且解向量长度不少于变量数，才进入统一后验证。
\fld{proven\_optimal} 还要求：
\begin{enumerate}
  \item 规范化状态字符串属于实现识别的最优状态集合；
  \item MIP gap 有限；
  \item $gap\le\fld{mip\_gap}+10^{-9}$。
\end{enumerate}
识别集合包括 \texttt{optimal}、\texttt{highs optimal}、
\texttt{optimal solution found}、\texttt{optimality gap reached} 等固定字符串。一个后端若返回语义等价
但未列入的状态，仍会得到可行结果，但不会标记已证明最优。

\subsection{统一后验可行性检查}

所有启发式/外部/Native 解必须通过：
\begin{align}
\|A_{eq}x-b_{eq}\|_\infty&\le10^{-6},\\
\|(Ax-b)_+\|_\infty&\le10^{-6},\\
\operatorname{dist}(x_j,\{0,1\})&\le10^{-4}\quad(j\in\mathcal B),\\
l_i-10^{-6}\le x_i&\le u_i+10^{-6}.
\end{align}
任何一项失败都把 \fld{feasible} 改为 false 并返回；残差数值只在 verbose 日志中出现，没有存入
\texttt{TopoReconfResult}。通过后 \fld{feasible=true}，且
\fld{optimal=proven\_optimal}。

\subsection{目标值与分解}

模型内开关项对初始闭合边使用 $-\lambda_{sw}c_e\beta_e$，省略常数
$\lambda_{sw}c_e$；初始断开边使用 $+\lambda_{sw}c_e\beta_e$。因此
\fld{milp\_objective}$=c^Tx$ 可以为负，且不是四个正向业务成本的直接总和。

\texttt{ObjBreakdown} 采用更易读的正向动作计费：
\begin{align}
J_{switch}^{report}&=\lambda_{sw}\sum_e c_e[\beta_e\ne\beta_e^0],\\
J_{loss}^{report}&=\lambda_{loss}\sum_{e\in E_{line}}|r_e|
\begin{cases}t_e,&\fld{loss\_aware}\land\fld{enable\_pf},\\\beta_e,&\text{其他},\end{cases}\\
J_{shed}^{report}&=\lambda_{shed}\left(\sum_i s_i^P+\sum_{i\in AC}s_i^Q\right),\\
J_{island}^{report}&=\lambda_{island}\sum_{g\ne g_0}\gamma_g.
\end{align}
因为 $J_{switch}^{report}$ 恢复了省略常数，通常有
\begin{equation}
\fld{milp\_objective}\ne
J_{switch}^{report}+J_{loss}^{report}+J_{shed}^{report}+J_{island}^{report}.
\end{equation}
调用方不得用四项求和反向校验原始目标，除非同时恢复所有常数。

\subsection{损耗字段的三种不同口径}

\begin{center}\small
\begin{tabular}{@{}L{4.0cm}L{5.1cm}L{5.4cm}@{}}
\toprule
字段/数值 & 计算 & 含义\\
\midrule
\fld{obj\_terms.loss} & $\lambda_{loss}\sum |r|t$ 或 $\lambda_{loss}\sum|r|\beta$ & 加权目标贡献，无物理 MW 承诺\\
\fld{reconf\_loss\_mw} & $S_B\sum_{AC/DC\ closed}r_e$ & 假设每条闭边 1 pu 电流的拓扑代理\\
HTTP \texttt{reconfig\_loss\_mw} & 完整 PF 支路 $P_{from}+P_{to}$ 的正损耗和 & PF 收敛时的后验证损耗\\
\bottomrule
\end{tabular}
\end{center}

核心 \fld{base\_loss\_mw}、\fld{loss\_reduction\_mw}、\fld{loss\_reduction\_pct} 当前没有赋值
路径，保持默认 0；HTTP 使用局部变量另行计算并写入 JSON。\texttt{ONRResult::estimated\_loss\_mw}
映射自核心 \fld{reconf\_loss\_mw}，仍是拓扑代理。

\subsection{完整字段赋值真值表}

\begin{longtable}{@{}L{4.2cm}L{4.0cm}L{6.4cm}@{}}
\toprule
字段 & 核心赋值时机 & 解释/限制\\
\midrule
\endfirsthead
\toprule
字段 & 核心赋值时机 & 解释/限制\\
\midrule
\endhead
\fld{feasible} & 后验线性检查通过 & 不是 PF/OPF 可行\\
\fld{optimal}/\fld{proven\_optimal} & 状态+有限 gap 通过 & 启发式永远 false\\
结构化 open/closed/switched & 可行解提取 & 混合系统权威索引空间\\
旧式整数向量 & 同时填充 & 可能跨域碰撞\\
\fld{switch\_operations} & 状态发生变化 & 设备或分支稳定 ID\\
\fld{ordered\_switch\_actions} & 动作回投影后 & 静态操作顺序，不是时域仿真\\
\fld{total\_shed\_mw/mvar} & PF 变量存在时 & $s^P/s^Q$ 求和；无节点向量\\
\fld{milp\_objective} & $c^Tx$ & 含权重和省略常数\\
\fld{obj\_terms} & 可行解后处理 & switching 项与原目标常数口径不同\\
\fld{reconf\_loss\_mw} & 可行解后处理 & 名义电流代理\\
\fld{base\_loss\_mw} 等 & 不赋值 & 核心始终默认 0\\
\fld{solve\_time\_s} & 正常/大多数失败返回 & 早期空系统/无源返回仍为 0\\
\fld{solver\_backend/status/gap} & 求解路径 & 早期无源只有 status\\
\fld{bc\_stats} & 仅 Native 路径 & 外部后端保持默认\\
\fld{model\_scope} & 入口开始 & 即使后续不可行仍存在\\
\fld{post\_action.../post\_power.../full\_hybrid.../executable} & 不赋值 & 仅 HTTP 层修改副本\\
\bottomrule
\end{longtable}

\subsection{ValidityFlags 的准确读法}

核心入口在求解前即设置部分标志：\fld{radial\_topology\_enforced}、
\fld{device\_capability\_constraints\_enforced}、\fld{protection\_interlocks\_enforced} 为 true；
\fld{ac\_lindistflow\_enforced} 等于 \fld{enable\_pf}；DC/VSC 标志按资产存在设置。
这些标志描述“模型构造意图”，并不以 \fld{feasible} 为前提。

特别地，\fld{allow\_dc\_mesh=true} 时 DC 边不受树边数约束，但
\fld{radial\_topology\_enforced} 仍为 true。调用方必须联合读取选项/运行时拓扑结果，不能把该单个
布尔值解释为 AC、DC、VSC 全域都已证明径向。

\subsection{失败模式}

\begin{itemize}
  \item 空母线或无 AC/DC 线路：返回默认不可行结果；只有 VSC 而无线段也在此边界内；
  \item 无可用根源：返回不可行并给出明确 \fld{solver\_status}；
  \item 后端无 incumbent：返回不可行并保留后端失败摘要；
  \item 后验残差、整数性或边界失败：返回不可行，详细残差只在 verbose 日志；
  \item 无效故障/候选稳定 ID：混合核心多数静默忽略；历史 ONR 旧死代码中的 throw 不可达；
  \item 未知 \fld{solver}：落入 auto 行为，不拒绝输入。
\end{itemize}

% compatibility_api.tex —— AC 拓扑查询与历史 ONR 包装

\section{拓扑查询与历史 ONR 兼容接口}\label{sec:nr-compat}

\subsection{AC 拓扑查询}

三个查询只读取 \texttt{ACSystem.buses} 和在运 \texttt{ACBranch}：
\begin{align}
\srcpath{count\_islands(ac)}&=\text{DFS connected-component count},\\
\srcpath{is\_connected(ac)}&\Longleftrightarrow count=1,\\
\srcpath{is\_radial(ac)}&\Longleftrightarrow |E_{on}|=|V|-1\land is\_connected.
\end{align}
空系统的约定为 \texttt{count\_islands=0}、\texttt{is\_connected=true}、
\texttt{is\_radial=true}；单母线无支路也为 connected/radial。

端点不在母线表中的在运支路会被 DFS 邻接构造跳过，但仍被
\srcpath{count\_in\_service\_branches} 计入边数。因此无效端点可能使 \texttt{is\_radial} 失败；
查询不返回诊断。它们不考虑 DC 支路、VSC、Switch/CircuitBreaker 富组件或 projection。

\subsection{ONROptions 到 TopoReconfOptions 的映射}

\srcpath{solve\_optimal\_reconfiguration(const ACSystem\&)} 当前是兼容包装。映射固定为：
\begin{center}\small
\begin{tabular}{@{}L{5.0cm}L{5.0cm}L{4.5cm}@{}}
\toprule
ONR 输入 & 核心选项 & 固定/差异\\
\midrule
\fld{v\_min\_pu/v\_max\_pu} & 同名 & 直接复制\\
\fld{default\_rate\_mva} & 同名 & 直接复制\\
\fld{big\_m} & \fld{big\_m\_v} & 直接复制\\
\fld{max\_time\_s/mip\_gap/verbose} & 同名 & 直接复制\\
\fld{switchable\_branch\_ids} & 同名 & 空时展开为全部 AC 支路 ID\\
-- & \fld{enable\_pf/voltage/thermal} & 固定 true\\
-- & \fld{loss\_aware} & 固定 true\\
-- & $\lambda_{sw},\lambda_{loss},\lambda_{shed},\lambda_{isl}$ & $0,1,10^4,10^5$\\
-- & \fld{solver} & 固定 \texttt{highs}，实现可再 fallback SCIP\\
\bottomrule
\end{tabular}
\end{center}

因此 ONR 头文件仍描述的旧变量布局 $[\alpha,P,Q,v,f,t]$ 与自有 B\&C/MIR 不是当前可达路径。
实际模型包含混合核心的根变量、源聚合、切负荷和设备语义；纯 AC 包装只是在输入中没有 DC/VSC。

\subsection{ONRResult 映射}

核心结构化结果只筛选 \texttt{AC\_Line} 类别写入
\fld{open\_branch\_ids}/\fld{closed\_branch\_ids}。其他字段映射为
\begin{align}
\fld{feasible}_{onr}&=\fld{feasible}_{topo},\\
\fld{optimal}_{onr}&=\fld{optimal}_{topo}\lor\fld{proven\_optimal}_{topo},\\
\fld{milp\_objective}_{onr}&=\fld{milp\_objective}_{topo},\\
\fld{estimated\_loss\_mw}_{onr}&=\fld{reconf\_loss\_mw}_{topo},\\
\fld{bc\_stats}_{onr}&=\fld{bc\_stats}_{topo}.
\end{align}
由于包装固定外部 HiGHS 首选，\fld{bc\_stats} 通常是默认值而非实际 HiGHS 统计。

核心可行时，包装把 closed ID 集应用到 ACSystem 副本，再调用
\srcpath{solve\_power\_flow}（\fld{max\_iter}=200、\fld{tol}=$10^{-6}$）填充
\fld{verification\_pf}。PF 不收敛不会把 \fld{feasible} 改为 false，调用方必须单独检查
\fld{verification\_pf.converged}。

\subsection{兼容失败回退}

若核心返回不可行，但输入 AC 拓扑本来 connected，包装会在原始拓扑上运行 PF。若该 PF 收敛，
返回值被改写为：
\begin{center}
\fld{feasible=true},\quad \fld{optimal=false},\quad
\fld{milp\_objective=0},
\end{center}
open/closed 恢复为原始状态，\fld{estimated\_loss\_mw} 取原闭合支路
$S_B\sum r_e$。\texttt{ONRResult} 没有字段标明“这是未重构的原拓扑 PF 回退”，也不保留核心失败
原因。故 \fld{feasible=true} 在此兼容 API 中可能只代表基准拓扑 PF 可行，不代表 ONR MILP 得到
候选解。

\begin{quote}\small\textbf{重要：}
需要严格优化证书、混合域引用、切负荷或回退透明度的新代码，应直接调用
\srcpath{run\_topology\_reconfiguration}。历史 ONR API 只适合保留旧调用形状，并必须联合检查
\fld{optimal} 与 \fld{verification\_pf.converged}。
\end{quote}

\subsection{HybridPowerSystem 重载}

混合重载先调用 \srcpath{project\_to\_canonical\_models(sys)}，然后只把 \texttt{proj.ac} 传给
ACSystem 兼容入口。因此 Transformer、Switch、FlexibleLoad 等可经投影进入规范 AC 表，但 DC 网络、
DC 负荷、VSC 和 DC 故障语义均被丢弃。该重载不是混合 AC/DC ONR；混合系统应调用
\srcpath{run\_topology\_reconfiguration}。

\subsection{summary() 口径}

\texttt{ONRResult::summary()} 打印可行/最优、目标、代理损耗、AC 整数支路 ID 和 BCStats；
\texttt{TopoReconfResult::summary()} 打印目标、动作计数、\fld{reconf\_loss\_mw}、耗时以及旧式整数
列表。两个摘要都不会打印 \fld{model\_scope}、ValidityFlags、结构化类别、切负荷、后端状态或
已知限制，不能直接作为审计报告。

% runtime_http_contract.tex —— 生产 HTTP 路由同步契约

\section{HTTP/JSON 与 GUI 同步契约}\label{sec:nr-http}

\subsection{生产路由与会话语义}

生产入口为
\begin{center}
\texttt{POST /api/session/run\_reconfig}
\end{center}
实现于 \srcpath{tests/run\_gui\_server.cpp}。它读取当前会话的富模型快照，执行\textbf{单时刻}
重构，并为了沿用前端时序表格把结果包装为 \fld{num\_steps=1}。响应中的
\fld{steps\_topology}、\fld{sw\_on\_per\_step} 等不是多时段优化证据。

\subsection{请求选项}

路由从顶层 JSON 或 \texttt{options} 对象读取：
\begin{itemize}
  \item 电压：\texttt{v\_min\_pu}/\texttt{v\_max\_pu}，兼容别名
        \texttt{v\_min}/\texttt{v\_max}；
  \item 求解：\texttt{mip\_gap}、\texttt{max\_time\_s}（HTTP 默认 60 s）、
        \texttt{solver}、\texttt{skip\_heuristic}；
  \item 模型：\texttt{enable\_pf}、\texttt{enable\_voltage}、
        \texttt{enable\_thermal}、\texttt{split\_domain\_trees}、
        \texttt{allow\_dc\_mesh}、\texttt{loss\_aware}（HTTP 默认 true）；
  \item 目标：四个 \texttt{lambda\_*} 与 \texttt{max\_switch\_ops}；
  \item 故障：\texttt{line\_failures}（AC 兼容）、\texttt{fault\_ac}、
        \texttt{fault\_dc}、\texttt{fault\_vsc}。
\end{itemize}
路由不接收显式 \texttt{switchable\_branches}；候选由富模型设备类型、角色、能力和锁状态自动生成。
当系统同时含 DC 母线与 VSC 且请求未显式给 \texttt{split\_domain\_trees} 时，路由自动启用分域模式。

HTTP 默认值与 C++ 默认值存在两处有意差异：\fld{max\_time\_s} 为 60 而非 300，
\fld{loss\_aware} 为 true 而非 false。复现实验必须记录调用层。

\subsection{后动作验证流水线}

核心 \fld{feasible} 后，路由执行：
\begin{enumerate}
  \item 将请求故障和 \fld{switch\_operations} 写入富模型副本；
  \item 重新 canonical projection，\fld{strip\_dead\_islands=false}；
  \item 分别检查 AC、DC 和混合图的连通/孤岛/径向状态；
  \item 在富模型上运行 \srcpath{solve\_power\_flow}；
  \item 运行 \srcpath{solve\_ac\_opf}；异常只转为验证消息；
  \item 仅当动作映射、序列、回投影、PF、OPF 全部通过时设置
        \fld{executable=true}。
\end{enumerate}

路由中的 \fld{full\_hybrid\_opf\_validated} 实际由 \srcpath{solve\_ac\_opf(rich\_reconfig)} 的
\fld{converged} 赋值。其具体混合覆盖取决于 AC OPF 门面与输入资产，不应仅凭字段名推断所有
DC/VSC 非线性约束都得到独立认证。

\subsection{响应字段分组}

\begin{longtable}{@{}L{4.4cm}L{4.2cm}L{6.0cm}@{}}
\toprule
组 & 代表字段 & 口径\\
\midrule
\endfirsthead
\toprule
组 & 代表字段 & 口径\\
\midrule
\endhead
核心候选 & \srcpath{feasible}、\srcpath{milp\_feasible}、\srcpath{milp\_objective}、\srcpath{obj\_terms} & 核心线性模型\\
工程可执行 & \srcpath{executable}、\srcpath{zero\_load\_loss\_valid}、验证消息 & PF/OPF/动作联合\\
模型范围 & \texttt{model\_scope}, \texttt{validity} & 核心+HTTP 后处理标志\\
求解器 & \texttt{solver\_name}, \texttt{solver\_status} & 后端或启发式；当前不返回 gap\\
切负荷 & \srcpath{total\_shed\_mw}、\srcpath{total\_shed\_mvar} & 核心 LinDistFlow 变量总和\\
真实 PF 损耗 & \srcpath{base\_loss\_mw}、\srcpath{reconfig\_loss\_mw}、reduction & 仅 PF 收敛时有物理意义\\
拓扑 & base/reconfig 的 radial、connected、islands 字段 & AC/DC/混合分组\\
AC 兼容表 & open/closed AC ID 与 \srcpath{branch\_details} & 只列 AC\\
混合详情 & \srcpath{dc\_branch\_details}、\srcpath{vsc\_details} & 带 domain 的独立数组\\
设备动作 & \srcpath{switch\_operations}、\srcpath{ordered\_switch\_actions} & 富设备/分支稳定 ID\\
伪时序视图 & \texttt{num\_steps=1}, \texttt{steps\_topology} & 单快照包装\\
\bottomrule
\end{longtable}

\fld{zero\_load\_loss\_valid} 定义为
\begin{equation}
\fld{executable}\land \fld{total\_shed\_mw}\le10^{-6},
\end{equation}
因此核心候选切负荷为零但后验证失败时，该字段仍为 false。这是比单看
\fld{total\_shed\_mw} 更严格的 GUI 结论。

\subsection{已审计的损耗响应缺陷}

当前路由写入
\begin{equation}
\texttt{estimated\_loss\_mw}
=\fld{milp\_objective}\times\fld{base\_mva}.
\end{equation}
该式量纲不成立：\fld{milp\_objective} 可能同时含加权开关、切负荷、孤岛和损耗项，且开关项省略
常数。它既不同于核心 \fld{reconf\_loss\_mw}，也不同于 HTTP 的 PF
\texttt{reconfig\_loss\_mw}。在代码修复前，客户端必须忽略 \texttt{estimated\_loss\_mw}，
只在 \texttt{reconfig\_pf\_converged=true} 时使用 \texttt{reconfig\_loss\_mw}。

\subsection{失败与 HTTP 状态}

核心不可行通常仍返回 HTTP 200，\texttt{milp\_feasible=false}；JSON 解析、输入访问或求解过程中
抛出的 \texttt{std::exception} 返回 400；未知异常返回 500。HTTP 200 不等于拓扑可行，更不等于
\fld{executable=true}。

\subsection{GUI 消费约束}

GUI 应：
\begin{itemize}
  \item 用 domain+ID 展示混合边，AC 旧式 ID 列表只用于纯 AC 表格；
  \item 同时展示 MILP 可行、PF/OPF 后验证和可执行性，不把它们折叠为一个“成功”；
  \item 把 \fld{milp\_objective} 标为无量纲加权目标，把损耗标为 PF 后验证值；
  \item 在 \fld{proven\_optimal} 未由 HTTP 暴露的现状下，不显示“全局最优”；
  \item 把单时刻响应标为一次重构，而非调度曲线。
\end{itemize}

% verification_audit.tex —— 验证矩阵、审计结论与边界

\section{验证矩阵与系统审计}\label{sec:nr-verification}

\subsection{审计范围}

本轮逐段核对：
\begin{itemize}
  \item 两个公共头的全部类型、默认值、注释与入口；
  \item \srcpath{topology\_reconfiguration.cpp} 的全部 1949 行可达核心路径；
  \item \srcpath{topology\_analysis.cpp} 的查询、兼容包装和包装后不可达旧模型；
  \item CMake 源文件接线、五个直接/交叉测试目标；
  \item \srcpath{POST /api/session/run\_reconfig} 的请求映射、富模型回投影、PF/OPF 与 JSON；
  \item web/e2e 消费点、理论速查、模块文档地图和审计台账。
\end{itemize}
“Deep”表示公共接口、核心实现、失败/结果语义和注册测试均已检查；不表示所有网络规模、所有求解器
组合或所有富设备类型都已动态穷举。

\subsection{注册测试映射}

\begin{longtable}{@{}L{4.1cm}L{6.1cm}L{4.4cm}@{}}
\toprule
目标 & 主要覆盖 & 本轮 Release 实测\\
\midrule
\endfirsthead
\toprule
目标 & 主要覆盖 & 本轮 Release 实测\\
\midrule
\endhead
\srcpath{test\_topology\_crossval} &
AC 查询、ExternalGrid/DC 源、AC/DC 负荷聚合、无源失败、候选稳定 ID、同号 AC/DC、DC 故障、
旧 ONR、6-bus 枚举、33-bus PF、DER & 12 cases / 86 assertions\\
\srcpath{test\_reconfig\_options} &
分域树、双 VSC、DC mesh、G4/G5、loss aware、Fuse、可操作 tie、上游保护三步序列、动作预算 &
10 cases / 37 assertions\\
\srcpath{test\_device\_flows} &
设备终端潮流与 dead-section/tie 投影保持 & 5 cases / 18 assertions\\
\srcpath{test\_distribution\_pipeline} &
rich model $\to$ projection $\to$ PF $\to$ OPF $\to$ SC $\to$ ONR $\to$ JSON &
1 case / 178 assertions\\
\srcpath{test\_crossmodule\_integration} &
IEEE 33-bus ONR $\to$ DC OPF $\to$ Newton PF $\to$ carbon & 2 cases / 82 assertions\\
\bottomrule
\end{longtable}

上述五个 \texttt{macos-release} 目标均由当前源码重新构建后运行并返回 0，共 30 cases、
401 assertions。这是 focused Release rebuild，不替代完整 1600+ 项回归、sanitizer 或生产 HTTP
E2E；逐项数值和依赖提交见第~\ref{sec:nr-numerical-crossval} 节。

\subsection{覆盖充分的契约}

已有直接回归锁定：
\begin{itemize}
  \item AC/DC 同号 ID 使用 \texttt{BranchRef} 消歧；旧 AC 整数候选不会误操作 DC；
  \item ExternalGrid-only、DC-only 源与无源失败；
  \item AC bus demand + Load、DCBus demand + DCLoad 的加法聚合；
  \item 分域 AC/DC 树、VSC bridge 自由、DC mesh 和 DC 独立开环；
  \item Fuse 不参与即时恢复、tie 直接闭合、隔离设备上游 trip/operate/reclose；
  \item 6-bus 与穷举、33-bus 与后续 PF/配电 PF、跨模块 OPF/PF/碳流；
  \item verbose=false 不向 stdout/stderr 泄漏求解器诊断。
\end{itemize}

\subsection{动态覆盖缺口}

下列高风险契约尚无直接、专门的断言：
\begin{itemize}
  \item \texttt{solver=auto/highs/scip/native} 的完整分派、fallback 与状态/gap 规范化；
  \item 图启发式成功路径的 \fld{proven\_optimal=false} 和无限 gap；
  \item 后端返回成功但后验等式/不等式/整数性/变量界失败的拒绝分支；
  \item 负权重、$v_{min}>v_{max}$、非正容量、未知 solver、未匹配故障 ID 的输入拒绝；
  \item 切负荷不超过节点需求、P/Q 同比例削减和纯无功/负负荷边界；
  \item CircuitBreaker 的锁定/操作能力与 Switch 规则的一致性；
  \item 核心 \fld{base\_loss\_mw} 等默认字段、HTTP \texttt{estimated\_loss\_mw} 口径；
  \item ONR 核心失败后“原拓扑 PF 可行”回退的可观察标志；
  \item HTTP \fld{executable} 的注册生产路由专项 E2E；现有参数库脚本未注册到 CMake。
\end{itemize}

\subsection{开放审计发现}

\subsubsection{NR-01：HTTP estimated loss 量纲错误（高）}

证据：\srcpath{run\_gui\_server.cpp} 的 \texttt{POST /api/session/run\_reconfig} 把
\fld{milp\_objective}$\times S_B$ 写入 \texttt{estimated\_loss\_mw}。目标含多种权重和省略常数，
不能转换为 MW。影响是 GUI/API 客户端可能发布错误损耗值。修复应删除该字段、改名为目标缩放诊断，
或映射到明确的核心代理；真实损耗只在 PF 收敛后报告。需加 HTTP 回归区分切负荷/开关成本与损耗。

\subsubsection{NR-02：ValidityFlags 在求解前宣称约束已执行（中）}

证据：核心在空系统/无源/后端失败之前即把 radial/device/protection 标志设为 true；
\fld{allow\_dc\_mesh=true} 时仍标 radial=true。影响是不可行结果或 DC mesh 可能被误读为全域已认证。
修复应把“requested/modelled/enforced/validated”拆开，至少在不可行返回时不声明 validated，并增加
AC/DC 分域字段。

\subsubsection{NR-03：solver 公共注释与实际 fallback 漂移（中）}

证据：\texttt{TopoReconfOptions::solver} 注释声称 auto 为 HiGHS$\to$Native，实际为
HiGHS$\to$SCIP；显式 highs 也会 fallback SCIP。\texttt{topology\_analysis.hpp} 仍声称现行 ONR
直接调用自有 B\&C/MIR，而包装在旧模型前已返回。影响是部署方可能错误估计依赖、可复现性和证书来源。
修复应同步头注释，决定是否恢复 Native fallback，并对四种字符串注册分派测试。

\subsubsection{NR-04：核心结果存在永不赋值的公共字段（中）}

\fld{base\_loss\_mw}、\fld{loss\_reduction\_mw}、\fld{loss\_reduction\_pct} 和四个后验证/
可执行性字段在核心无赋值路径。HTTP 在局部副本/局部变量上补值，但直接 C++ 调用只见默认值且无法
区分“未评估”与“评估失败”。修复应移除字段、改为 optional/availability，或把后验证作为显式 API。

\subsubsection{NR-05：历史 ONR 回退隐藏 MILP 失败（中）}

核心不可行但原 AC 拓扑 PF 收敛时，兼容包装返回 \fld{feasible=true}、目标 0、原拓扑支路集，且没有
fallback/limitation/status 字段。调用方可能把它当作 ONR incumbent。修复应保留核心失败原因并新增
\texttt{fallback\_used/model\_scope}，或不改变 \fld{feasible} 的优化语义。

\subsubsection{NR-06：旧 AC MILP 为大段不可达实现（低）}

\srcpath{topology\_analysis.cpp} 在兼容包装的无条件 return 后仍保留约 600 行旧模型。它增加审计成本，
且头注释、测试名称容易继续引用旧行为。修复应删除不可达代码并把历史模型移入文档/版本历史，或恢复
为明确、独立且可测试的入口；不应继续让一份函数体表达两个互斥实现。

\subsection{已知模型边界}

\begin{itemize}
  \item 单时刻、稳态、正序/平衡等值；无三相相域拓扑变量；
  \item LinDistFlow 不含支路损耗，P/Q 热限为矩形，VSC 有功无效率损耗且 Q 只接 AC 端；
  \item \fld{loss\_aware} 是 $r|P|$ 线性代理，不是 $I^2R$；默认非 loss-aware 更只是 $r\beta$；
  \item 根与源容量是聚合近似，ExternalGrid 能力以短路容量或默认容量替代运行限额；
  \item 切负荷无需求上界与恒功率因数联动；
  \item 设备能力只对投影可归因的 AC Switch/CircuitBreaker 较完整，DC/VSC 结构化候选是分支级；
  \item 保护联锁是静态元数据规则，不含短路电流、保护配合、同步或操作时延；
  \item 结果不包含逐母线电压/切负荷、逐边 P/Q、对偶量、下界、约束残差和模型限制列表；
  \item 核心不执行完整 PF/OPF，兼容 ONR 只验证 AC PF，HTTP 才追加 PF/AC OPF；
  \item 旧整数 ID 结果不适用于 AC/DC/VSC 同号系统。
\end{itemize}

\subsection{工业使用判据}

把结果用于运行建议前，至少应满足：
\begin{enumerate}
  \item 输入通过严格 validation，故障与候选稳定 ID 全部匹配；
  \item 使用 \texttt{BranchRef}，并保存实际 solver/backend/status/gap；
  \item \fld{feasible=true}；若宣称最优，还须 \fld{proven\_optimal=true}；
  \item 动作已回投影、\fld{switch\_sequence\_valid=true}；
  \item 富模型重投影后 AC/DC 拓扑符合业务要求；
  \item 完整 PF 收敛并满足电压/热限，所需场景下完整 OPF 也收敛；
  \item 损耗来自收敛 PF，而非目标或名义电阻代理；
  \item 对关键故障做备用动作、保护配合和现场可操作性复核。
\end{enumerate}

% numerical_cross_validation.tex -- 重构数值结果与独立交叉验证
\section{数值结果与独立交叉验证}\label{sec:nr-numerical-crossval}

\subsection{复现实验条件}

本节针对“只有断言、没有数值结果”的缺口，给出可从日志复核的结果和口径。实验
使用 Apple Silicon macOS/arm64 的 \srcpath{macos-release}（Release、\texttt{-O3}、
\texttt{-DNDEBUG}），主仓库提交为 \texttt{b1bf15736c96}，
本地 \texttt{../MIPSolvers} 为 \texttt{7721936245d5}，
与仓库声明 pin \texttt{3bf1e66749e3} 不一致。该依赖差异必须随结果一起报告。

本轮 focused Release rebuild 目标为：
\begin{verbatim}
cmake --build --preset macos-release --target \
  test_topology_crossval test_reconfig_options test_device_flows \
  test_distribution_pipeline test_crossmodule_integration -j4
\end{verbatim}
五个目标共通过 30 cases/401 assertions。完整 CTest、HTTP 生产 E2E、sanitizer
和“pinned dependency”重跑均未在本轮执行。

\subsection{6-bus：MILP 与穷举生成树}

对 6-bus 小系统枚举所有 spanning trees，独立最优目标为
0.04290000 pu；HiGHS MILP 返回 0.0429 pu，按
\((J_{\rm MILP}-J_{\rm exhaustive})/J_{\rm exhaustive}\) 计算相对差为 0\%。
测试将此比较标为 informational，且结果字段 \fld{optimal=false}；因此只能说
当前 incumbent 与穷举最优值数值一致（门槛 $\le10^{-8}$），不能写成 MILP 已证明
最优。MILP 选中的闭合支路数为 5，恰为 \(n_{bus}-1\)，满足该实例的径向性检查。

加入 DER（renewable + storage）后，MILP 目标从无 DER 的 0.0429 pu 降至
0.026 pu，降幅约 39.3939\%。独立 BFS/DPF 路径给出的 DER 方案损耗为
0.0381 MW；这是物理潮流损耗的另一个模型输出，不是 0.026 pu MILP 目标的单位
换算，也不能用它证明 MILP 目标代表损耗。

\subsection{IEEE 33-bus：ONR 输出与独立 BFS}

IEEE 33-bus ONR 返回线性损耗 proxy 0.075197（无量纲标幺目标），开断支路为
10、13、15、28、33，闭合支路数为 32，等于 \(33-1\)。日志中的
\texttt{estimated loss = 14.531 MW} 是生产路由把目标乘以 base MVA 的已知量纲错误，
本手册明确禁止把它当作物理损耗。该运行还报告 \texttt{B\&C nodes explored=0}、
\texttt{gap=1e+30}，所以没有 proven-optimal 证据。

将相同开关拓扑交给独立 BFS/DPF loss evaluator，基准拓扑损耗为 0.118446 MW，
重构拓扑为 0.0995629 MW，下降约 15.9424\%。这组数值证明“重构后该独立潮流模型
的损耗下降”，但不与 MILP proxy 进行单位混算，也不替代 Newton PF 的完整电压/热限
证书。

\subsection{IEEE 33-bus 跨模块流水线}

\texttt{test\_crossmodule\_integration} 以同一 case 串联 ONR、DC OPF、Newton PF
和碳流追踪。基准 Newton PF residual 为 \(7.4675\times10^{-9}\)，基准物理有功
损耗 0.202677 MW；重构后流水线 PF residual 为 \(2.50789\times10^{-9}\)，损耗
0.16527 MW，按 PF 口径下降约 18.4565\%。碳流节点平衡误差为
\(3.16036\times10^{-7}\%\)。DC OPF 基准与重构目标均为 74.3，说明该实例的
OPF 目标没有因拓扑动作改变，但不表示 PF 损耗不变。

这里必须保留两套损耗交叉验证：
\[
\begin{aligned}
\text{BFS/DPF:}\quad &0.118446\longrightarrow0.0995629\ {\rm MW},\\
\text{Newton PF:}\quad &0.202677\longrightarrow0.16527\ {\rm MW}.
\end{aligned}
\]
前者来自独立 BFS/DPF，后者来自 Newton PF/跨模块流水线；二者算法、网络模型或
参数口径不同，不能合并成一个“唯一真实损耗”。

\subsection{选项、设备与保护动作的数值检查}

\texttt{test\_reconfig\_options -s} 的 HiGHS 状态为
\texttt{StrictHiGHS Optimal run=0}。分域检查中 DC mesh 的 \fld{dc\_open}=0，
AC/DC split-domain radiality 返回 \fld{dc\_closed}=1、\fld{ac\_closed}=1。
loss-aware 实例保留必要 tie（闭合支路数 2），设备/保护路径验证动作序列
\([20\ \mathrm{open},21\ \mathrm{close},20\ \mathrm{reclose}]\)，并返回
\fld{switch\_sequence\_valid=true}。这些是选项和动作约束的证据，不是完整保护配合
或现场同步检验。

\subsection{结论、门槛与限制}

数值准入规则为：6-bus MILP 与穷举相对差 $\le10^{-8}$ 且仍检查
\fld{optimal=false}；33-bus BFS 必须收敛并通过 radiality；跨模块 PF 必须独立报告
residual，不能以 MILP objective 代替。当前证据通过这些门槛，但以下结论仍未闭环：
\begin{itemize}
  \item 33-bus 当前运行没有 proven optimality，不能发布全局最优声明；
  \item \texttt{estimated\_loss\_mw} 仍是错误字段，真实损耗只能取收敛 PF/BFS 的明确口径；
  \item 未重跑完整 HTTP E2E、sanitizer 和 pinned MIPSolvers；
  \item 核心 MILP 仍是单时刻线性/名义损耗模型，不覆盖非线性 AC OPF、保护配合和操作时延。
\end{itemize}


\section{跨模块工业级技术评价基线}

本章规定本手册所述模块进入工程算例、批量研究或生产服务前必须共同满足的评价基线。
具体模型公式、变量、约束和专用算法以正文相应章节为准；本章不扩大源码已经实现的范围。

\subsection{输入、输出、边界条件与数据结构审计}
输入审计应逐项检查有限性、单位、上下界、枚举值和引用完整性。系统对象必须先按所需层级完成
校验；AC 与 DC 母线采用域限定映射，组件稳定 \texttt{index}、向量位置和临时图索引不得混用。
输出审计必须记录单位、索引空间、模型范围、收敛或可行状态、后端、容差、迭代/节点计数和告警。
空系统、空时域、孤岛、重复 ID、零容量、退化约束与不可用可选后端均属于边界测试集。

\subsection{工程假设与数值稳定性分析}
模型使用的平衡、线性化、稳态、独立性、概率分布、时间离散或网络等值假设必须与应用场景匹配。
数值评价至少包含变量缩放、绝对/相对残差、可行性违反、条件数或枢轴质量、步长/阻尼、迭代停滞
和随机种子复现性。若采用正则化、平滑、回退、松弛或近似，应同时报告触发条件、参数值、对主指标
的敏感性以及是否改变模型口径；不得用“已返回结果”替代收敛或有效性证明。

\subsection{复杂度与资源分析}
令 $n_x$、$n_c$、$n_t$、$n_s$、$|V|$、$|E|$ 分别表示变量、约束、时段、场景、节点和边数。
报告应分别给出建模、求解、后处理的时间与内存。图遍历通常为 $O(|V|+|E|)$；逐时段/逐场景
流程至少线性增长为 $O(n_t n_s)$ 次子问题调用；稠密线性代数最坏可达 $O(n_x^3)$，稀疏方法则应
报告结构非零数、填充率和因子分解峰值内存；MILP/NLP 不给出虚假的多项式最坏界，而报告节点数、
间隙、时限和实测扩展曲线。复杂度结论必须基于固定硬件、固定后端和可复现算例族。

\subsection{异常工况处理}
非法输入应在进入求解器前返回结构化错误；不可行、无界、不收敛、时限、内存不足和后端缺失必须
相互区分。部分结果只有在对应 \fld{ValidityFlags}、\fld{model\_scope}、
\fld{model\_limitations} 或等价字段允许时才可使用。异常路径不得静默丢弃 AC/DC 资产、制造平衡
节点、把 NaN 改写为零或把近似解标为全局最优；系统原始对象应保持可恢复和可重试。

\subsection{与其他模块的接口关系}
接口链路遵循“工程语义模型 $\to$ 校验 $\to$ 投影/装配 $\to$ 求解或分析 $\to$ 结果回投影”。
跨模块传递时保留稳定 ID、单位、符号、时间基准和场景身份；消费潮流、OPF、动态或可靠性结果前，
先检查其收敛、覆盖范围和有效性标志。HTTP/JSON/Python 层只做语义保持适配，不重新解释求解结果。

\subsection{测试与验证建议}
最低验证矩阵包括：手算小例、标准算例、边界/异常输入、随机性质测试、算法或后端交叉校核、
序列化往返、稳定 ID 回投影、AC/DC 同号母线、重复运行确定性和规模扩展测试。数值算法还应加入
残差独立重算、雅可比有限差分或自动微分校核；近似模型应与高保真模型比较并给出误差分布，而非
只报告单个平均值。回归报告必须列明构建配置、算例版本、容差、通过/失败/跳过数及未验证范围。

\subsection{工业级评估指标}
\begin{itemize}
  \item \textbf{正确性}：守恒残差、约束最大违反、解析/基准误差、结果回投影一致性；
  \item \textbf{鲁棒性}：收敛率、不可行识别率、异常分类准确率、扰动敏感性与随机复现率；
  \item \textbf{性能}：端到端时延、吞吐量、峰值内存、并行效率与规模增长斜率；
  \item \textbf{可追溯性}：输入模式版本、稳定 ID 覆盖率、模型范围/限制字段完整率、告警可定位率；
  \item \textbf{工程有效性}：与标准、外部引擎或现场数据的偏差，以及误判/漏判对业务指标的影响。
\end{itemize}
指标应预先给出验收阈值和失败处置；未达到阈值时只能标记为实验性或受限适用。

\subsection{适用范围、局限性与后续改进}
本手册只评价当前源码覆盖的模型、后端和数据结构，不对未建模物理过程、未安装求解器或未执行的
外部验证作保证。后续改进应按“缺口 $\to$ 可量化指标 $\to$ 源码/测试变更 $\to$ 独立验证”闭环，
优先消除会造成错误结论的语义缺口，其次提升数值鲁棒性与性能，最后扩展模型范围。


\end{document}
